\chapter{Introducción, motivación y objetivos}
En este capítulo se presentan la motivación, los objetivos y la organización de la memoria del proyecto TransDrakkar, realizado en la asignatura Procesos de Ingeniería del Software.
\section{Motivación}
La gestión de usuarios permite abordar de forma práctica problemas habituales de un producto software: registro, autenticación, autorización, persistencia y despliegue. El proyecto permite trabajar estos aspectos de manera incremental y relacionar las decisiones técnicas con la planificación, las pruebas y la revisión de resultados.
La elaboración de esta memoria también sirve para aprender a documentar el proceso y utilizar Overleaf y LaTeX de cara a futuros trabajos académicos, incluido el TFG.
\section{Objetivos}
El objetivo principal es desarrollar y documentar un prototipo web de transporte de mobiliario con gestión de usuarios aplicando un proceso de desarrollo incremental.
Objetivos específicos:
\begin{itemize}
    \item Aprender a usar una metodología para desarrollar un producto software, registrando la planificación y el seguimiento en DoneTonic.
    \item Desarrollar un prototipo por capas con registro, acceso, sesiones y permisos diferenciados.
    \item Incorporar persistencia de cuentas y preparar el despliegue de la aplicación.
    \item Integrar confirmación por correo y acceso con Google, completando su validación real antes de considerarlos terminados.
    \item Conectar solicitud del cliente, revisión administrativa, planificación de rutas y entrega por conductor, con seguimiento e incidencias.
    \item Evaluar los resultados mediante pruebas y evidencias del incremento obtenido.
\end{itemize}
\section{Estructura de la memoria}
El primer capítulo presenta la motivación y los objetivos. El segundo queda reservado para antecedentes opcionales. El tercero describe la propuesta, la metodología y las herramientas. El cuarto recoge los requisitos, la arquitectura, la evolución del proyecto, el uso de la aplicación, las pruebas y el despliegue. El quinto presenta conclusiones del cierre del Sprint 2 y trabajo futuro. El anexo de desarrollo reúne el detalle de los sprints y las evidencias de DoneTonic y su revisión.
